\section{Week 3}

\subsection{Assignments}
\begin{itemize}
    \item Exercises, 7/10 approval required
    \item Helicopter lab
    \item Matlab assessments
\end{itemize}


\subsection{Grading}
\begin{itemize}
    \item Final exam: 80\%
    \item Matlab assessments: 20\%
\end{itemize}


\subsection{Course material}
\begin{itemize}
    \item Numerical Optimization, J.Nocedal and S. J Wright, 2nd ed.
    \item Errata on BB.
    \item Note on Merging Optimization and Control, B. Foss and T.A. Heirung.
    \item Note on Matrix Calculus, T.A. Heirung.
\end{itemize}

\subsection{Motivation}
Optimization: Find best solution to a given task (mathematical problem).
For the non general problems numerical techniques must be used.
Applied mathematics built on:
\begin{itemize}
    \item Convexity and non-convexity
    \item Global vs. local solution
    \item Constrained vs. unconstrained problems
    \item Feasible set
\end{itemize}

Typical application for Optimization:
\begin{itemize}
    \item Autonomous vehicles
    \item Power plants
    \item Oil rigs
    \item Machine learning
    \item Finance
\end{itemize}

Minimization and Maximization, convention this course is for minimization
Maximizing f(x) will be seen as minimizing -f(x).
\begin{equation}
        \min_{x \in \mathbb{R}} f(x)
\end{equation}

What characterizes an optimum?
\begin{itemize}
\item \[\min_{x \in \mathbb{R^n}} f(x)\]
\item \[f \in \mathbb{C{^2}}\]
\item \[ \frac{d}{dx} f(x*) = 0\]
\item \[ \frac{d^2}{dx^2} f(x*) > 0\]
\end{itemize}

To make use of optimization, must define an objective which is a quantitative
measure of the performance of a system:
\begin{itemize}
    \item profit
    \item time
    \item energy
    \item or any specific quantity
\end{itemize}
Objectives depends on variables, and the goal is to find values  of the variables that 
optimize the objective. The variables are often constrained or bounded by some way.

Modeling defines formulating an objective, variables and constraints for a given taks.
There after an optimization algorithm can be defined, for example in machine learning minimizing MSE.

\begin{itemize}
    \item x is the vector of variables
    \item f is the objective function of x returning a scalar which we want to minimize
    \item $c_i$ are constraint functions of x returning a scalar and defines equations or inequalities that the unknown vector x must satisfy.
\end{itemize}


\subsection{(Un)Constrained Optimization}

Constraints are mathematical inequalities for the variables $x_n$ for the desired task.

Unconstrained:
\begin{itemize}
    \item \[\min_{x \in \mathbb{R^n}} f(x)\]
\end{itemize}

Constrained:
\begin{itemize}
    \item \[\min_{x \in \mathbb{R^n}} f(x) s.t. C_i(x) = 0, i \in \mathbb{\epsilon} C_i(x)>=0, i \in \mathbb{I}\].
    \item \[C_i(x)\] constrained function
    \item $\epsilon$: Equality constraint
    \item $I$: Inequality constraint
\end{itemize}

Feasible set:
\begin{itemize}
    \item \[\Omega = {x \in \mathbb{R^n} s.t. C_i(x) = } \]
\end{itemize}

\subsection{Linear programming}
\begin{equation}
    \min_{x \in \mathbb{R} c^Tx s.t.}
        \left\{ \begin{aligned}
            a_i^T x = b_i , i \in \epsilon \\
            a_i^Tx \geq b_i, i \in I
        \end{aligned}
        \right.
\end{equation}

\subsubsection{Farmer example, find  max profit or min neg profit}

\subsection{Quadratic programming}
\begin{equation}
    \min{x \frac{1}{2}x^TGx + c^Tx s.t. \begin{aligned}
        a_i^T x = b_i , i \in \epsilon \\
        a_i^Tx >= b_i, i \in I \\
        \end{aligned}}
\end{equation}

\subsection{Convexity}
Convex set: A set is convex if
\begin{equation}
    x,y \in \mathbb{S} if \alpha x + (1-\alpha)y \in \mathbb{S}, \forall \alpha \in [0,1]
\end{equation}

Convex function: f:S -> R is a convex function if S is convex and
$every x,y \in \mathbb{S} : f(\alpha x + (1-\alpha)y <= \alpha f(x) + (1-\alpha)f(y) , \forall \alpha \in [0,1])$
"A line from point x to y should be below the function from f(x) to f(y)

An optimization problem is convex if
 \begin{itemize}
    \item f(x) is a convex function
    \item The feasible set is a convex set
 \end{itemize}

\subsubsection{Local and global solutions}
x* is a global solution if $f(x*) <= f(x) all x \in\Omega$.
Let N be a neighborhood of x:
x* is a local solution if $f(x*) <= f(x) \forall x \in N \subseteq \Omega$

Omega is the set of all x values which fullfill the constraints.